\documentclass[12pt,letterpaper]{article}



%%
%% Required packages:
%%
\usepackage{etex}
\usepackage{amsmath}
\usepackage{amssymb}
\usepackage{amstext}
\usepackage{amsthm}
\usepackage{fancyhdr,fancybox}
\usepackage{mathrsfs} % need for \mathscr command
\usepackage{multicol}
\usepackage[normalem]{ulem}
%\usepackage{pictex,pictexplus}
% \usepackage{pictexwd}
\usepackage{rotating}
\usepackage{url}
\usepackage{xfrac}
\usepackage{booktabs}
\usepackage{color,soul}
\usepackage{comment}

\usepackage{hyperref}
\hypersetup{
    colorlinks=true,     % false: boxed links; true: colored links
    linkcolor=blue,      % color of internal links
    citecolor=blue,      % color of links to bibliography
    filecolor=blue,      % color of file links
    urlcolor=blue        % color of external links
}


\usepackage{tikz}
\usetikzlibrary{backgrounds,calc}

%%
%% Common formatting commands
%%
\setlength{\parindent}{0in}
\setlength{\textwidth}{7in}
\setlength{\evensidemargin}{-0.25in}
\setlength{\oddsidemargin}{-0.25in}
\setlength{\parskip}{.5\baselineskip}
\setlength{\topmargin}{-0.5in}
\setlength{\textheight}{9in}

%               Problem and Part
\newcounter{problemnumber}
\newcounter{partnumber}
\newcounter{subpartnumber}
\newcommand{\Problem}{%
  \refstepcounter{problemnumber}%
  \setcounter{partnumber}{0}%
  \setcounter{subpartnumber}{0}%
  \item[\makebox{\hfil\textbf{\theproblemnumber.}\hfil}]%
  }
\newcommand{\InProblem}[2][1.6]{%
  \refstepcounter{problemnumber}%
  \setcounter{partnumber}{0}%
  \setcounter{subpartnumber}{0}%
  \textbf{\theproblemnumber.}\ \ \parbox[t]{#1in}{#2}%
}
\newcommand{\InSmallProblem}[1]{\InProblem[1.05]{#1}}
\newcommand{\NoProblem}[1]{%
  \mbox{\hfil\textbf{#1}\hfil}%
  }
\newcommand{\UnProblem}{%
  \refstepcounter{problemnumber}%
  \setcounter{partnumber}{0}%
  \setcounter{subpartnumber}{0}%
  \mbox{\hfil\textbf{\theproblemnumber.}\hfil}%
  }
\newcommand{\InFlexProblem}[1]{%
  \refstepcounter{problemnumber}%
  \setcounter{partnumber}{0}%
  \setcounter{subpartnumber}{0}%
  \textbf{\theproblemnumber.}\ \ #1%
}
%%  Part:
\newcommand{\Part}{%
  \renewcommand{\thepartnumber}{(\alph{partnumber})}%
  \refstepcounter{partnumber}
  \setcounter{subpartnumber}{0}%
  \item[(\alph{partnumber})]%
}
\newcommand{\InPart}[2][1.75]{%
  \renewcommand{\thepartnumber}{(\alph{partnumber})}%
  \refstepcounter{partnumber}%
  \setcounter{subpartnumber}{0}%
  (\alph{partnumber})\ \ \parbox[t]{#1in}{#2}%
}
\newcommand{\UnPart}{%
  \refstepcounter{partnumber}%
  \renewcommand{\thepartnumber}{(\alph{partnumber})}%
  \setcounter{subpartnumber}{0}%
  (\alph{partnumber})%
}
\newcommand{\InLargePart}[1]{\InPart[3.05]{#1}}
\newcommand{\InFullPart}[1]{\InPart[6.2]{#1}}
\newcommand{\InSmallPart}[1]{\InPart[1.05]{#1}}
%% SubPart:
\newcommand{\SubPart}{%
  \refstepcounter{subpartnumber}%
  \item[(\roman{subpartnumber})]%
  }
\newcommand{\InSubPart}[2][2.25]{%
  \stepcounter{subpartnumber}%
  (\roman{subpartnumber})\ \ \parbox[t]{#1in}{#2}%
}
\newcommand{\InSmallSubPart}[1]{\InSubPart[1.5]{#1}}
\newcommand{\InTinySubPart}[1]{%
  \stepcounter{subpartnumber}%
  (\roman{subpartnumber})\ \ \makebox{#1}%
}
\newcommand{\InMySubPart}[2][2.25]{%
  \stepcounter{subpartnumber}%
  \parbox[t]{#1in}{\makebox[0.3in][l]{(\roman{subpartnumber})}\ \ {#2}}%
}
\newcommand{\TrueFalse}{\textbf{T}\ \ \textbf{F}\ \ }
\newcommand{\TFPart}[1]{% 
  \stepcounter{partnumber}%
  \setcounter{subpartnumber}{0}%
  \item[(\alph{partnumber})] \TrueFalse\parbox[t]{5.5in}{#1}}

\newcommand{\Points}[1]{(#1 points) \ }
\newcommand{\Point}[1]{(#1 point) \ }


%% For Answers:
\newcounter{answernumber}
\newcommand{\Answer}{%
  \refstepcounter{answernumber}%
  \setcounter{partnumber}{0}%
  \setcounter{subpartnumber}{0}%
  \item[\makebox{\hfil\textbf{\theanswernumber.}\hfil}]%
  }


% Setting up the headers for the first page
%\chead{} % will default to what is typed in the file.
\fancyhead[L]{\textsc{Math 22a}}
\fancyhead[R]{\textsc{Fall 2024}}
\fancyfoot[R]{
	\vspace*{\fill}\footnotesize\textcopyright{} \emph{Harvard Math 22a}	
}
\renewcommand{\headrulewidth}{0pt}

%\fancyhead[C]{}
%	\chead{}	
%	\lhead{}
%	\rhead{}
%	\cfoot{}


% Putting copyright on the footer of each page
%\fancypagestyle{secondpage}{
%	\fancyhead[C]{TESTing again} % change this to be blank
%		%\chead{}	
%	\fancyhead[L]{bears} % change this to be blank
%	\fancyhead[R]{dogs} % change this to be blank
%	\fancyfoot[R]{ %keeps the footer the same
%		\vspace*{\fill}\footnotesize\textcopyright{} \emph{Harvard Math 22a}	
%	}
%}
%\renewcommand{\headrulewidth}{0pt}
%\pagestyle{fancy}
\pagestyle{empty}


%\lhead{}
%\rhead{}
%\rfoot{\vspace*{\fill}\footnotesize\textcopyright{} \emph{Harvard Math 22a}}
%\renewcommand{\headrulewidth}{0pt}
%\pagestyle{fancy}




%%
%% Common shortcut commands:
%%
\newcommand{\ds}{\displaystyle}
\newcommand{\sgn}{\operatorname{sgn}}
\renewcommand{\Pr}{\operatorname{P}}
\newcommand{\CPr}[3][{}]{\Pr_{\text{#1}}\left(#2\,|\,#3\right)}

\newcommand{\C}{\mathbb{C}}
\newcommand{\R}{\mathbb{R}}
\newcommand{\Z}{\mathbb{Z}}
\newcommand{\N}{\mathbb{N}}
\newcommand{\Q}{\mathbb{Q}}
\newcommand{\D}{\mathbb{D}}

\newcommand{\rref}{\operatorname{rref}}
\newcommand{\im}{\operatorname{im}}
\newcommand{\rank}{\operatorname{rank}}
\newcommand{\diag}{\operatorname{diag}}
\newcommand{\Span}{\operatorname{span}}
%\renewcommand{\vec}[1]{\mathbf{#1}}
\newcommand{\charpoly}{\mathscr{P}}
\newcommand{\nicefrac}[2]{\sfrac{#1}{#2}} % using xfrac package
\newcommand{\topology}{\mathcal{T}}
\newcommand{\Inn}{\operatorname{Inn}}

\newcommand{\blankbox}[2]{\fbox{\rule{#1}{0in}\rule{0in}{#2}}}

\newenvironment{myindentpar}[1]%
 {\begin{list}{}%
         {\setlength{\leftmargin}{#1}
          \setlength{\rightmargin}{\leftmargin}}%
         \item[]%
 }
 {\end{list}}
\newtheorem{theorem}{Theorem}
\newtheorem*{NStheorem}{Theorem}
\newtheorem{cor}[theorem]{Corollary}
\newtheorem*{claim}{Claim}
\newtheorem{defn}{Definition}

\newenvironment{amatrix}[1]{%
  \left(\begin{array}{@{}*{#1}{c}|c@{}}
}{%
  \end{array}\right)
}

%--------grstep
% For denoting a Gauss' reduction step.
% Use as: \grstep{\rho_1+\rho_3} or \grstep[2\rho_5 \\ 3\rho_6]{\rho_1+\rho_3}
\newcommand{\grstep}[2][\relax]{%
   \ensuremath{\mathrel{
       {\mathop{\longrightarrow}\limits^{#2\mathstrut}_{
                                     \begin{subarray}{l} #1 \end{subarray}}}}}}
\newcommand{\swap}{\leftrightarrow}


\usetikzlibrary{arrows}
\usepackage{wrapfig}

\makeatletter
\renewcommand*\env@matrix[1][*\c@MaxMatrixCols c]{%
	\hskip -\arraycolsep
	\let\@ifnextchar\new@ifnextchar
	\array{#1}}
\makeatother

\def\env@matrix{\hskip -\arraycolsep
	\let\@ifnextchar\new@ifnextchar
	\array{*\c@MaxMatrixCols c}}

\chead{\Large\textbf{Coordinates, $\vec\S$4.4}}% 

\begin{document}
\thispagestyle{fancy}

Recall from previously:

\begin{itemize}
\item 	The subset $\{ v_1, v_2, \dots, v_p\}$ of the vector space $V$ is {\bf linearly independent} if $$c_1 v_1+\dots+ c_p v_p =0_V$$ has only the trivial solution.
\item {\bf Theorem:} Given $\{ v_1, \dots, v_p\}$ with $p\geq 2$ and $v_1\neq 0_V$. The vectors $v_1, \dots, v_p$ are linearly dependent if and only if for some $j>1$, the vector $v_j$ is a linear combination of $v_1, \dots, v_{j-1}$.
\item Let $H$ be a subspace of a vector space $V$. The set $\{v_1, \dots, v_p\}$ is a {\bf basis} for $H$ if
	\begin{enumerate}
		\item  $\{v_1, \dots, v_p\}$ is linearly independent,
		\item $H=\text{Span}\{v_1, \dots, v_p \}$.
	\end{enumerate}
\item {\bf Spanning Set Theorem:} Let $S=\{ v_1,\dots,v_p\} \subseteq V$ and $H=\text{Span}\{v_1,\dots, v_p \}$.
\begin{enumerate}
	\item Suppose $v_k$ can be written as a linear combination of the other vectors in $S$, then $S-\{v_k\}$ still spans $H$.
	\item If $H\neq \{ 0_V\}$, some subset of $S$ is a basis for $H$.
\end{enumerate} 
\item {\bf Theorem:} The pivot columns of an $m \times n$ matrix $A$ form a basis for $\text{Col}(A)$.
\end{itemize}

\newpage

{\bf Theorem:} Let $\mathcal{B} =\{b_1, \dots, b_n\}$ be a basis for a  vector space $V$. Then for each $x \in V$ there exists unique scalars $c_1,\dots, c_n$ such that $x=c_1 b_1+\dots+ c_n b_n$.



\begin{enumerate}


\Problem Prove the theorem.

\vfill

\begin{defn} Suppose $\mathcal{B}=\{b_1, \dots, b_n\}$ is a basis for a vector space $V$. The {\bf coordinates of $x$ relative to $\mathcal{B}$} are the weights $c_1, \dots, c_n$ such that $x=c_1 b_1 +\dots c_n b_n$. We define the {\bf coordinate mapping} from $V$ to $\mathbb{R}^n$ by $[x]_\mathcal{B}=\begin{bmatrix} c_1\\ \vdots\\ c_n \end{bmatrix}$.
\end{defn}

\Problem Let $V=\mathbb{R}^3$ with basis $\mathcal{E}=\{e_1, e_2, e_3\}$. Let $x= 3e_1 +4e_2+5e_3$, then find $[x]_\mathcal{E}$. Let $\mathcal{B}=\left\{\begin{bmatrix} 1\\ 0\\ 0\\\end{bmatrix},\begin{bmatrix} 1\\ 1\\ 0\\\end{bmatrix},\begin{bmatrix} 1\\ 1\\ 1\\\end{bmatrix} \right\}.$ Find $[x]_\mathcal{B}$.

\vfill

\newpage

\Problem Let $b_1=\begin{bmatrix} 1\\ 1\\ \end{bmatrix}$ and $b_2=\begin{bmatrix} 2\\ 1\\ \end{bmatrix}$, then $\mathcal{B}=\{b_1, b_2\}$ is a basis for $\mathbb{R}^2$. Let $x =\begin{bmatrix} 1\\ 0\\ \end{bmatrix}$. Find $[x]_\mathcal{B}$.

\vfill

\begin{defn} Let $\mathcal{B}=\{b_1, \dots, b_n\}$, then define $P_\mathcal{B}=[b_1 \; \dots \; b_n]$. We call $P_\mathcal{B}$ the {\bf change of coordinates matrix}.
\end{defn}
Note $x = P_\mathcal{B} [x]_{\mathcal{B}}$ and $[x]_\mathcal{B}=P_{\mathcal{B}}^{-1} x.$

{\bf Theorem:} Let $\mathcal{B}=\{ b_1, \dots, b_n\}$ be a basis for a vector space $V$. Then the coordinate mapping $x \mapsto [x]_\mathcal{B}$ is a one-to-one, onto, linear transformation from $V$ to $\mathbb{R}^n$.

\Problem Prove the theorem.

\vfill

\newpage

\begin{defn}
	Let $T: V \to W$ be a linear map of vector spaces. If $T$ is one-to-one and onto, then $T$ is a {\bf vector space isomorphism}, and we say $V$ and $W$ are isomorphic. We write $V \cong W$.
\end{defn}

\Problem Show that $\mathbb{P}_3$ and $\mathbb{R}^4$ are isomorphic.

\vfill

{\bf Theorem:} If a vector space $V$ has a basis $\mathcal{B}=\{b_1, \dots, b_n\}$, then any set in $V$ containing more than $n$ vectors is linearly dependent.

\Problem Prove the theorem.

\vfill

{\bf Cor:} If $V$ is a vector space with basis $\mathcal{B}=\{b_1, \dots, b_n\}$, then each set of linearly independent vectors has at most $n$ elements.

{\bf Theorem:} If a vector space $V$ has a basis with $n$ elements, then every basis of $V$ must have exactly $n$ vectors.

\Problem Prove the theorem.

\vfill

\begin{defn}
If $V$ is spanned by a finite set, then $V$ is called {\bf finite dimensional,} and the {\bf dimension} of $V$, written $\text{dim}(V)$, is the number of vectors in a basis of $V$. Define $\text{dim}(\{0_V\})=0$. If $V$ is not spanned by any finite set, then $V$ is called {\bf infinite dimensional}.
\end{defn}

\end{enumerate}

\begin{comment}
\newpage
\chead{\Large\textbf{Coordinates -- Answers / Solutions}}%
\lhead{}
\rhead{}
\thispagestyle{fancy}

\begin{enumerate}
  \Answer 
  
  \begin{proof}
  Since $\mathcal{B}$ is a basis, we know there exists scalars $c_1, \dots, c_n \in \mathbb{R}$ such that $$x=c_1 b_1 +\dots +c_n b_n.$$ In order to prove the scalars are unique, we suppose $d_1, \dots, d_n$ are scalars such that $x=d_1 b_1+\dots +d_n b_n$. Now subtracting the two representations of $x$, we get$$0=x-x= (c_1 b_1 +\dots +c_n b_n) -(d_1 b_1+\dots +d_n b_n).$$ Using the vector space axioms, we rearrange the terms and factor to get $$0= (c_1-d_1)b_1+\dots+ (c_n -d_n)b_n.$$ Since $\mathcal{B}$ is a basis, we know $b_1, \dots, b_n$ are linearly independent, and hence $$c_1-d_1=0, \dots, c_n-d_n=0.$$ Thus $c_1=d_1,\dots, c_n=d_n$. Thus the scalars are unique.
  \end{proof}

\Answer $[x]_\mathcal{E}=\begin{bmatrix} 3\\ 4\\ 5\\ \end{bmatrix}$. In order to find $[x]_\mathcal{B}$ we want to find scalars $c_1, c_2, c_3$ such that $x= c_1 b_1 +c_2 b_2+c_3 b_3$. Thus we form the augmented matrix $[b_1 \; b_2 \; b_3 \; | x]$ and solve
$$\begin{amatrix}{3}
1 & 1 & 1 & 3\\
0 & 1 & 1 & 4\\
0 & 0 & 1 & 5\\
\end{amatrix}\sim
\begin{amatrix}{3}
1 & 0 & 0 & -1\\
0 & 1 & 0 & -1\\
0 & 0 & 1 & 5\\
\end{amatrix}.$$ Thus $x=(-1)b_1 + (-1) b_2+5b_3$, and hence $[x]_\mathcal{B}=\begin{bmatrix} -1\\ -1\\ 5\\ \end{bmatrix}$.

\Answer Again we form the augmented matrix corresponding to the vector equation $x=c_1 b_1 +c_2 b_2$:
$\begin{amatrix}{2}
1 & 2 & 1\\
1 & 1 & 0\\
\end{amatrix}\sim\begin{amatrix}{2}
1 & 0 & -1\\
0 & 1 & 1\\
\end{amatrix}$. Thus $[x]_\mathcal{B}=\begin{bmatrix} -1\\ 1\\ \end{bmatrix}$.

\Answer

\begin{proof} Let $\mathcal{B}=\{b_1,\dots, b_n\}$ be a basis for the vector space $V$. Let $T: V \to \mathbb{R}^n$ be given by $T(x)=[x]_\mathcal{B}$.

{\bf Subclaim 1:} $T$ is linear.

Let $x,y\in V$. Since $\mathcal{B}$ is a basis, there exists (unique) scalars $c_1, \dots, c_n, d_1, \dots, d_n$ such that $$x=c_1 b_1+\dots+ c_n b_n,$$ and $$y=d_1 b_1+\dots+ d_n b_n.$$
Thus using the vector space axioms, we get
$$
[x+y]_\mathcal{B}=[(c_1+d_1)b_1+\dots+(c_n+d_n)b_n]_\mathcal{B}=\begin{bmatrix} c_1+d_1\\ \vdots\\ c_n+d_n\\ \end{bmatrix}=\begin{bmatrix} c_1\\ \vdots\\ c_n\\ \end{bmatrix}+\begin{bmatrix} d_1\\ \vdots\\ d_n\\ \end{bmatrix}=[x]_\mathcal{B}+[y]_\mathcal{B}.$$

Let $c\in \mathbb{R}$, then $[cx]_\mathcal{B}=[cc_1b_1+\dots +c c_n b_n]_\mathcal{B}=\begin{bmatrix} c c_1\\ \vdots\\ c c_n\\ \end{bmatrix}=c [x]_\mathcal{B}$.

Thus $T$ is linear.

{\bf Subclaim 2:} $T$ is one-to-one (injective).

Let $x, y \in V$ such that $T(x)=T(y)$. By the first theorem today, the coordinates are unique, so $x=y$.

{\bf Subclaim 3:} $T$ is onto (surjective).

Let $f=\begin{bmatrix} f_1 \\ \vdots \\ f_n\end{bmatrix} \in \mathbb{R}^n$, then $f_1b_1+\dots +f_n b_n \in V$ and $T(f_1b_1+\dots +f_n b_n)= f$. Thus $T$ is surjective.
\end{proof}
\Answer Let $\mathcal{B}=\{1, t, t^2, t^3 \}$ be the standard basis on $\mathbb{P}_3$. By the previous theorem the coordinate mapping with respect to basis $\mathcal{B}$ is a vector space isomorphism.

\Answer 

\begin{proof}
Let $a_1,\dots, a_p \in V$ with $p>n$, then $[a_1]_\mathcal{B}, \dots, [a_p]_\mathcal{B} \in \mathbb{R}^n$. We have more than $n$ vectors in $\mathbb{R}^n$, so we proved they must be linearly dependent. Thus there exists $c_1,\dots, c_p \in \mathbb{R}$, not all zero, such that  $$0 = c_1 [a_1]_\mathcal{B}+\dots+c_p[a_p]_\mathcal{B}.$$ Since the coordinate mapping is linear, we get $$0=[c_1 a_1 + \dots + c_p a_p]_\mathcal{B}.$$ Since the coordinate mapping is injective, we get $0=c_1 a_1 + \dots +c_p a_p$ where not all $c_1, \dots, c_p$ are zero. Thus the vectors are linearly dependent in $V$.
\end{proof}

\Answer

\begin{proof}
Let $\mathcal{B} =\{b_1,\dots, b_n\}$ be a basis for $V$ and let $\mathcal{C} =\{c_1, \dots, c_m\}$ be a basis for $V$. Since $\mathcal{B}$ is a set of linearly independent vectors in $V$, the previous theorem tells us that $n \leq m$. Since $\mathcal{C}$ is a set of linearly independent vectors in $V$, the previous theorem tell us that $m \leq n$. Thus $n=m$. 
\end{proof}

\end{enumerate}  

\end{comment}
\end{document}


%%% Local ables: 
%%% mode: latex
%%% TeX-master: t
%%% End: 
